\bookchapter{第五十一章：介面開發——Jetpack Compose 入門}

\chapterintro{介面係用戶見到嘅所有嘢。喺 Compose 出現之前，安卓介面用 XML 寫——好似填空題咁麻煩；2021 年穀歌推出咗 Jetpack Compose：直接用 Kotlin 程式碼「畫」介面，所見即所得，改一行程式碼預覽即刻刷新。到 2026 年，Compose 已經係新項目嘅絕對主流。呢一章，我哋學識「畫畫」。}

\booksection{一、Compose 嘅思想：介面就係函數}

Compose 嘅核心思想一句講曬：介面 = 函數。每樣可以顯示喺螢幕上面嘅嘢（文字、按鈕、圖片）都係一個「可組合函數」（Composable），用 @Composable 標記。將佢哋好似積木噉嵌套埋一齊，就係介面。

\begin{lstlisting}[language=Kotlin]
@Composable
fun Greeting(name: String) {
    Text(text = "你好，$name！")     // 一嚿"文字積木"
}

// 巢狀組合：大積木套小積木
@Composable
fun HomeScreen() {
    Column {                          // 直向排列
        Greeting(name = "張三")
        Greeting(name = "李四")
        Button(onClick = { /* 㩒咗做乜 */ }) {
            Text("㩒我")
        }
    }
}
\end{lstlisting}

留意 @Composable 呢頂「帽」——佢話畀編譯器知「呢個係一嚿介面積木」。積木可以無限嵌套，呢個就係 Compose 嘅全部秘密。

\booksection{二、基礎組件：文字、按鈕、輸入框、圖片}

Compose 嘅組件就係「積木」，最常見嘅四種：

\begin{center}
\begin{tabularx}{\linewidth}{|X|X|X|}
\hline
\textbf{組件} & \textbf{做乜嘢} & \textbf{必備參數} \tabularnewline
\hline
Text & 顯示文字 & text（要顯示嘅字） \tabularnewline
Button & 按鈕 & onClick（點擊回調） \tabularnewline
TextField & 輸入框 & value、onValueChange（輸入變化回調） \tabularnewline
Image & 顯示圖片 & painter（圖片資源） \tabularnewline
\hline
\end{tabularx}
\end{center}

\begin{center}\small\color{gray}Compose 四大基礎組件\end{center}

\begin{lstlisting}[language=Kotlin]
@Composable
fun BasicWidgets() {
    Column {
        Text(text = "呢度係一段文字")
        Button(onClick = { println("被㩒咗") }) {
            Text("按鈕")
        }
        var input by remember { mutableStateOf("") }
        TextField(
            value = input,
            onValueChange = { input = it },
            label = { Text("輸入啲乜嘢") }
        )
    }
}
\end{lstlisting}

輸入框嗰兩行入面出現咗 remember \{ mutableStateOf(“”) \}——呢個就係 Compose 嘅「狀態」魔法，本章第五節會詳細講。而家先混個面熟。

\booksection{三、佈局三件套：Column、Row、Box}

組件有咗，點樣擺放？記住三個「盒」：

\begin{itemize}
  \item Column：打直排（好似排隊打飯噉）；
  \item Row：打橫排（好似並排坐噉）；
  \item Box：疊住放（好似啤牌疊埋一齊，適合做角標、遮罩）。
\end{itemize}

\begin{lstlisting}[language=Kotlin]
Column {
    Text("第一行")
    Row {
        Text("左邊")
        Text("右邊")
    }
    Box {
        Text("底層")
        Text("疊喺上面嘅角標")
    }
}
\end{lstlisting}

想佈局更加好睇？畀組件加 Modifier（修飾符）——佢就好似畀積木「上妝」噉：

\begin{lstlisting}[language=Kotlin]
Text(
    text = "靚嘅文字",
    modifier = Modifier
        .padding(16.dp)          // 四圍留白
        .fillMaxWidth()          // 撐滿闊度
        .background(Color.Blue)  // 背景色
        .clickable { /* 㩒 */ }// 㩒得
)
// 鏈式調用：一個接一個，次序有意義（先 padding 再背景，背景唔包邊距）
\end{lstlisting}

\begin{tcbwarn}{注意：Modifier 嘅次序有講究}
padding(16.dp).background(Blue) 同 background(Blue).padding(16.dp) 效果唔同——前者係「先留白再塗色」（背景唔包邊距），後者係「先塗色再留白」（背景包埋邊距）。新手最容易喺呢度畫錯「奇怪嘅白邊」。唔確定嗰陣，調轉次序試吓，預覽即刻話你知答案。

\end{tcbwarn}

\booksection{四、狀態：畀介面「活」起上嚟}

靜態介面冇乜意思。想撳掣之後數字 +1、輸入框可以打字，就需要狀態（State）。Compose 嘅規則：介面永遠忠實反映狀態，狀態一變，介面自動重畫。

\begin{lstlisting}[language=Kotlin]
@Composable
fun Counter() {
    // remember + mutableStateOf：記住呢個會變嘅值
    var count by remember { mutableStateOf(0) }

    Column {
        Text(text = "你㩒咗 $count 次")
        Button(onClick = { count++ }) {   // 改狀態
            Text("㩒我 +1")                // 界面自動刷新
        }
    }
}
\end{lstlisting}

三行程式碼，一個識郁嘅計數器。呢個就係 Compose 比舊式寫法正嘅地方：唔使手動「刷新介面」，改資料就得。記住兩個關鍵函數：remember（令個值喺介面重組嗰陣唔會被清空）同 mutableStateOf（話畀 Compose 知「呢個值變咗要重畫」）。

用戶點擊

狀態改變（count++）

Compose 感知到變化

自動重畫界面（重組）

\begin{center}\small\color{gray}圖 3：Compose 嘅「狀態 介面」單向資料流\end{center}

\booksection{五、主題同深色模式：一次配置，全局生效}

模板項目入面嘅 MyFirstAppTheme \{ ... \} 就係主題。佢基於 Material 3 設計規範，幫你將顏色、字體、形狀統一管理。支援深色模式只需要三步：

\begin{enumerate}
  \item 喺主題入面定義好淺色／深色兩套配色（模板已經自動生成咗 ColorScheme）；
  \item 系統轉去深色模式嗰陣，isSystemInDarkTheme() 會自動判斷；
  \item 用主題提供嘅顏色（例如 MaterialTheme.colorScheme.primary）而唔係寫死顏色值。
\end{enumerate}

\begin{tcbinfo}{小貼士：Material 3 動態取色}
Android 12 開始嘅「Material You」可以令 App 由你嘅桌布入面提取顏色。Compose 模板預設支援（dynamicColor）。你個 App 喺唔同用戶嘅手機上面會「自動換裝」，呢個就係現代安卓嘅浪漫。

\end{tcbinfo}

\booksection{六、預覽、動畫同實時編輯：Compose 嘅「外掛」}

Compose 三個開發體驗上嘅「外掛」，一用咗就返唔到轉頭：

\begin{itemize}
  \item @Preview 預覽：畀 Composable 加 @Preview 註解，Android Studio 右邊即刻 render 出個介面——改程式碼所見即所得，唔使每次都裝落手機度；
  \item 實時編輯（Live Edit）：2026 年嘅 Android Studio 支援喺模擬器／真機上面「熱更新」部分改動，改文字顏色呢類小調整，設備上面即刻變；
  \item 動畫 API：animateFloatAsState、AnimatedVisibility 幾行程式碼就可以令元素平滑過渡——「高級感」往往就差咗呢幾十行。
\end{itemize}

\begin{lstlisting}[language=Kotlin]
@Composable
fun AnimatedBox() {
    var big by remember { mutableStateOf(false) }
    // 值變嗰陣自動平滑過渡到目標值
    val size by animateDpAsState(if (big) 120.dp else 60.dp)

    Box(Modifier
        .size(size)
        .background(MaterialTheme.colorScheme.primary)
        .clickable { big = !big }
    )
}
\end{lstlisting}

\begin{tcbinfo}{小貼士：先功能後動畫}
新手成日犯嘅毛病係「介面都未整啱，就開始糾結動畫」。先令功能行通，最後先統一加動效——呢個亦都係「簡單至上」原則：動畫係錦上添花，功能係雪中送炭。

\end{tcbinfo}

\begin{tcbpractice}{嚟動手：整一個「心情按鈕」}
將模板項目改成一個心情切換器：一個按鈕，點擊之後喺 開心 難過 嬲 之間循環切換，文字同背景色跟住變。提示：用 when 根據狀態揀文案同顏色。做完傳畀朋友玩，收穫第一波「嘩你識整 App 喇」。

\end{tcbpractice}

\begin{tcbsummary}{本章小結}
\begin{itemize}
  \item Compose 思想：介面 = @Composable 函數，積木式嵌套
  \item 四大組件：Text / Button / TextField / Image
  \item 佈局三件套：Column 打直排、Row 打橫排、Box 疊放
  \item Modifier 鏈式「上妝」：padding、fillMaxWidth、background、clickable——次序有意義
  \item 狀態係靈魂：remember + mutableStateOf，改資料介面自動重畫
  \item 主題用 Material 3 + 動態取色，深色模式一次配置全局生效
\end{itemize}

\end{tcbsummary}
